% ch3 B.1 准粒子 (原书页 104–111 / PDF p119–p126)

\section{$N+1$ 体问题}

\subsection{准粒子}

让我们来谈谈我们开始这一番讨论时所举的那个特别简单的元激发例子，也就是半导体或绝缘体中的自由载流子。这个问题常常被称作"$N+1$ 体问题"，因为我们真正感兴趣的只是由添加这个额外粒子所带来的那些变化。元激发有两大普遍类型；上面这个例子属于其中被称为"准粒子"（quasi-particles）的那一类——之所以这样称呼，是因为它们非常接近于独立的 Hartree-Fock 单电子波函数中的粒子，而且与它们由之产生的介质中的个别粒子具有相同的对易规则和电荷——在通常情形下，它们是费米子。与之相对照的那一类元激发是集体激发（collective excitation），其中包括声子、等离激元、激子、自旋波等等。$\mathrm{He}_4$ 中的旋子则是一个反常的例子，兼具两种类型各自的性质。

准粒子激发 $q^{\dagger}_{k}$ 的特征在于：它与相应的单粒子近似 $c^{\dagger}_{k}$ 有一个有限的交叠：
\begin{equation*}
\left(q^{\dagger}_{k}\Psi_{\mathrm{g}},\ c^{\dagger}_{k}\Psi_{\mathrm{g}}\right) = Z^{1/2} \tag{1}
\end{equation*}
$Z$ 为有限值（不是 $1/N$ 的量级）。这必须被当作一个定义；我们无法证明这类激发总是存在，而且本来也不应当指望能证明，因为确实存在一些系统，在其中这些激发实际上并不存在。式 (1) 绝非平庸，因为我们当然知道，真实基态 $\Psi_{\mathrm{g}}$ 与 Hartree-Fock 近似 $\Psi_{\mathrm{HF}}$ 之间的交叠是 $e^{-N}$ 的量级。因此，在讨论准粒子时，我们所使用的概念只离开独立粒子模型有限的距离，这是一个巨大的概念上的有利之处。

对式 (1) 的一种思考方式——它同多体理论家实际处理准粒子的方式 (68) 相当贴近——是考虑这样的问题：假设在时刻 $t=0$，我用算符 $c^{\dagger}_{k}$ 作用于真实基态，这个算符（比如说）在一个平面波态 $e^{ik\cdot r}$ 中产生一个粒子。现在我等待一段很长的时间 $T$，在此期间哈密顿量作用于系统，于是 $T$ 时刻的新波函数是 $e^{-iHT}c^{\dagger}_{k}\Psi_{\mathrm{g}}$。现在我要问：粒子是否仍然处在态 $k$ 中？这意味着我施加湮灭算符 $c_{k}$，再与原来的基态作标量积。于是我做如下计算：
\begin{align*}
G_{k}(T) &= \left(e^{-iHT}\Psi_{\mathrm{g}},\ c_{k}\,e^{-iHT}c^{\dagger}_{k}\Psi_{\mathrm{g}}\right)\\
&= \left(\Psi_{\mathrm{g}},\ e^{iHT}c_{k}e^{-iHT}c^{\dagger}_{k}\Psi_{\mathrm{g}}\right)\\
&= \langle g|\ c_{k}(T)c^{\dagger}_{k}(0)\ |g\rangle \tag{2}
\end{align*}
现在，为了求出这个量，一个有益的做法是往这个矩阵元中插入一组严格的激发态，其中之一就是准粒子态 $q^{\dagger}_{k}\Psi_{\mathrm{g}}$；但正如我们说过的，这类态的数目可能非常之大。举例来说，$c^{\dagger}_{k}$ 很可能起到产生这样一个态的作用：它带有两个准粒子和一个准空穴（quasi-hole），三者的动量之和为 $k$：$q^{\dagger}_{k_{1}}q^{\dagger}_{k_{2}}q_{k_{1}+k_{2}-k}\Psi_{\mathrm{g}}$。

让我们把这类可能的激发态中随便哪一个记作 $\Psi_{n}$，并且注意：几乎所有可能类型的 $\Psi_{n}$ 都不是以有限的、而是以无穷小的振幅进入其中，因为具有同一动量的这类态有无穷多个。（唯一可能的例外是周期势中的其他能带态，这个例外处理起来很平凡。）于是
\begin{align*}
G_{k}(T) &= \sum_{m}\ \langle g|c_{k}|m\rangle\ e^{-i(E_{m}-E_{\mathrm{g}})T}\ \langle m|c^{\dagger}_{k}|g\rangle\\
&= Z e^{-iE'_{k}T} + \sum_{n}\ |\langle g|c_{k}|n\rangle|^{2}\ e^{-iE'_{n}T}
\end{align*}
这里已经把激发能 $E'$ 定义为 $E-E_{\mathrm{g}}$。在 $N\to\infty$ 的极限下，对 $n$ 的求和将变成如下形式的一个积分：
\begin{equation*}
\int e^{iE_{n}T}\ dE_{n}\ f(E_{n})
\end{equation*}
其中所有的 $E_{n}$ 都大于 $E_{k}$。这类积分总是导致急速衰减的函数；因此在长时间后的 $G(T)$ 就由那一个准粒子激发所支配。我们可以这样说：在态 $k$ 中产生这个电子之后，经过很长时间，它的 $Z$ 份额仍然留存下来，其余部分则已衰变而进入一个态的连续区；这一份量为 $Z$ 的剩余部分，其行为就好像它具有准粒子的能量一样。

注意，这为我们提供了一种"操作性的"（operational）定义准粒子算符 $q_{k}$ 的办法。也就是说，我们可以使用方程
\begin{align*}
\left(c^{\dagger}_{k}\Psi_{\mathrm{g}}\right)(T) &= e^{-iHT}c^{\dagger}_{k}\Psi_{\mathrm{g}}\\
&= Z^{1/2}e^{-i(E'_{k}+E_{\mathrm{g}})T}\ q^{\dagger}_{k}\Psi_{\mathrm{g}} + \int dE_{n}\ e^{-i(E'_{n}+E_{\mathrm{g}})T}\ f(E_{n})\Psi_{n}
\end{align*}
将此式乘以 $e^{i(E'_{k}+E_{\mathrm{g}})T}$ 并积分，得到
\begin{equation*}
q^{\dagger}_{k}\Psi_{\mathrm{g}}\ \propto \int_{0}^{\infty} dT\ e^{i(E'_{k}+E_{\mathrm{g}})T}\left(c^{\dagger}_{k}\Psi_{\mathrm{g}}\right)(T)
\end{equation*}
或者
\begin{equation*}
q^{\dagger}_{k} = \text{const.}\ \int_{0}^{\infty} dT\ e^{i(E'_{k}+E_{\mathrm{g}})T}\ c^{\dagger}_{k}
\end{equation*}
这就是说，我们只需对所有时间作适当的平均，就能从含时波函数中挑出属于真正准粒子的那一部分；这样做所给出的严格准粒子函数，比起任何背景成分来要多出无穷多倍。用能谱的语言来说，我们可以画出图 30。于是，如果乘上一个形如 $\int_{0}^{\infty}dT\ e^{-iET}$ 的能量 $\delta$ 函数滤波器，我们就能容易地把准粒子态及其能量挑出来。

%FIG{30}: 图 30

这种通过研究 $T\to\infty$ 极限下的行为来获得严格结果的技巧，是多体理论（M. B. T.）的一种特征性的形式工具。Kohn、Ambegaokar (69) 以及晚近的 Blount (45) 基于这些思想，已经就 $N+1$ 电子问题确立了大量严格结果。其中一个结果是严格地证明：在与单能带近似等价的理论中，可以把准粒子能量曲线 $E(k)$ 用作单能带哈密顿量中有效的动能部分，正如在 Bloch 电子理论中可以把哈密顿量处理成 $E(k)+V(R)$ 那样。另一个结果是证明：在长程极限下，势要被宏观的静态介电常数 $\kappa$ 所修正；当电子离电荷足够远时，它会以势 $eq/\kappa r$ 被一个静止电荷 $q$ 吸引。

这最后一个结果在文献中受到了极大的关注。虽然从物理观点来看它是平凡的，但数学上的论证却并非如此，而且它还揭示出一些有意思的东西。我在本课程中无法给出完整的推导，但我认为，粗略地展示一下其中牵涉到哪些东西，对于领会 $N+1$ 体问题的处理技巧是有教益的。

假设在材料中建立起了一个静电势 $V(r)$——比如说，通过引入外电荷 $e$——它当然会在哈密顿量中给出一个微扰项
\begin{equation*}
\int dr\ \Psi^{\dagger}(r)V(r)\Psi(r) = \frac{1}{N}\sum_{q,k}\ c^{\dagger}_{k+q}c_{k}\ V(q) \tag{3}
\end{equation*}
其中
\begin{equation*}
V(q) = \frac{1}{\Omega}\int dr\ e^{iq\cdot r}V(r) = \frac{4\pi e}{\Omega q^{2}}
\end{equation*}
这是对一个外电荷 $e$ 而言的。

我们在这里假定，引入电荷时离子实并不移动，因此势不会按相应于离子实运动的介电常数而折减；这在非极性半导体中是合理的，而在其他情形下这个效应无论如何也总可以被考虑进去。

式 (3) 取在两个准态（quasi-states）$K$ 与 $K+Q$ 之间的一个典型矩阵元是
\begin{equation*}
\left(\Psi_{K'},\ \sum_{k}c^{\dagger}_{k+Q}c_{k}\ \Psi_{K+Q}\right)V(Q) = M_{K}(Q)V(Q)
\end{equation*}
其中显然，一切 $q\neq Q$ 的矩阵元都必定因动量守恒而消失。Kohn 和 Ambegaokar 的定理是：该方程中的标量积等于 $1/\kappa$，这里 $\kappa$ 是电子气的介电常数——至少在 $Q\to 0$ 的极限下是如此，这一极限对应于非常长程的力［$V$ 的长程部分当然就是 $V(Q\to 0)$］。

这个定理乍看起来相当令人惊讶，因为恰恰在 $Q=0$ 的极限下，矩阵元 $\left(\Psi_{K},\ \sum_{k}n_{k}\ \Psi_{K}\right)$ 就等于粒子总数，即 $N+1$。

然而，出现这个分歧的原因是清楚的：单个准粒子所引起的扰动含有一些本质上是长程的成分，因为那个电子的电荷会把整个晶体一直极化到 $r\to\infty$，而这种长程效应在 $Q\to 0$ 的极限下导致奇特的行为。这确实是原因所在，而这一点可以用一种非常直接的办法来证明。让我们考察一个完整绝缘体的样品，往其中加入一个单独的电子——不是让它处在某个特定的准粒子态 $K$ 中，而是处在我早先提到过的那种波包（wave packet）中：它由围绕某一给定的 $K_{\mathrm{o}}$ 的那些态的线性组合构成，并且局域在一点 $R_{\mathrm{o}}$ 附近，我们将选取这一点作为原点：
\begin{equation*}
\Psi_{R_{\mathrm{o}}=0,\ K_{\mathrm{o}}} = \frac{\text{const.}}{\left[N(\Delta K)^{3}\right]^{1/2}}\ \sum_{K}\ e^{-\frac{1}{2}(K-K_{\mathrm{o}})^{2}/(\Delta K)^{2}}\ \Psi_{K} \tag{4}
\end{equation*}

让我们先从物理上看一看这个问题。我们已经在绝缘体的中心加入了一个额外的电子，它处在一个大小为 $\Delta R\approx 1/\Delta K$ 的区域内（图 31）。因此我们认为，在这个区域内有一份额外的电子电荷。现在，如果我们考虑由此造成的、从 $\Delta R$ 一直到样品外半径 $R$ 的那一层电介质球的极化，那么这个球将被电子电荷的电场 $E$ 所极化，极化强度为 $P=[(\kappa-1)/4\pi]E$。这就导致该球内、外表面上的总面电荷为
\begin{align*}
&+\ (\kappa-1)(\Delta R)^{2}E(\Delta R)\qquad\text{inner}\\
&-\ (\kappa-1)R^{2}E(R)\qquad\qquad\ \ \text{outer}
\end{align*}
而 $E$ 是由净电荷造成的场：
\begin{equation*}
E(\Delta R) = \frac{e-(\kappa-1)(\Delta R)^{2}E(\Delta R)}{(\Delta R)^{2}}
\end{equation*}
或者说 $Er^{2}=e/\kappa$，而面电荷为 $\pm[(\kappa-1)/\kappa]\ e$。

%FIG{31}: 图 31

实际上，球并不会在 $\Delta R$ 处的边界上被物理地切开，因此这个面电荷将近似地被内层介质球表面上的面电荷所补偿。极化电荷的实际位置只能粗略地指明为处在准粒子所在的地方。我们所能确切断定的只是：电子电荷中有 $(\kappa-1)/\kappa$ 的份额出现在球的外表面上，只有 $1/\kappa$ 的份额留存下来，从而能够与任何局域地插入的外来电荷 $q$ 相互作用。

这是所有涉及可极化介质中载流子的问题所共有的一种奇特现象的一个特别简单的例子：运动电荷中总有一个相当可观的份额实际上是处在样品表面上的。事实上，在金属中——那里 $\kappa(q\to 0)=\infty$——原则上全部电荷都在表面上；在电介质中则只有 $(\kappa-1)/\kappa$。换一种说法：那个无限长程的、$Q=0$ 的项总是表现出反常的、不连续的行为。我们也可以把它解释为一种屏蔽效应——至少部分地是如此；在长距离上，电子与其他电荷的相互作用势被削减了一个因子 $\kappa$。

还有待于指出这些事实与 Kohn-Ambegaokar 恒等式之间的关系，该恒等式涉及矩阵元
\begin{equation*}
\lim_{Q\to 0}M_{K}(Q) = \lim_{Q\to 0}\sum_{k}\left(\Psi_{K'},\ c^{\dagger}_{k+Q}c_{k}\ \Psi_{K+Q}\right) \tag{5}
\end{equation*}
注意，这个量如同常数 $Z^{1/2}$ 一样，是准粒子量与裸粒子（bare-particle）量之间的一个"交叉"标量积。事实上，它恰恰类似于量子场论中的"顶点重整化常数"（vertex renormalization constant），正如 $Z$ 之于该理论中的"波函数重整化"常数 $Z$。（译注：末句中前一个 $Z$ 原书即印作 $Z$，按上文的类比关系，疑应为 $Z^{1/2}$。）

我们通过计算一个包含 $\Delta R$ 球的球体内的电荷来求式 (5)，准粒子就处在 $\Delta R$ 球之内。引进一个光滑函数 $f(r)$，它在 $r=\Delta R$ 处等于 1，而在某个半径 $R_{1}$ 处下降到零，其中 $R\gg R_{1}\gg\Delta R$。于是，半径为 $R_{1}$ 的球体内的电荷由下式给出：
\begin{equation*}
e\left(\Psi_{R_{\mathrm{o}}=0,\ K_{\mathrm{o}}}\ \int dr\ f(r)\Psi^{\dagger}(r)\Psi(r)\ \Psi\right) \tag{6}
\end{equation*}
（译注：左端 bra 态的下标原书印作 $k_{\mathrm{o}}=0$；按式 (4) 及下文"$\Psi(R=0)$"的表述，应为 $R_{\mathrm{o}}=0$。）

让我们对这个表达式作 Fourier 分析，利用
\begin{equation*}
\Psi^{\dagger}(r)\Psi(r) = \sum_{Qk}\ e^{iQ\cdot r}\ c^{\dagger}_{k+Q}c_{k}
\end{equation*}
以及波包 $\Psi(R=0)$ 的 Fourier 展开式 (4)。式 (6) 成为
\begin{align*}
\frac{e\times\text{constants}}{(\Delta K)^{3}N}\ &\sum_{K,K',k,Q}\ \left(\int dr\ f(r)\,e^{iQ\cdot r}\right)\left(\Psi_{K'},\ c^{\dagger}_{k+Q}c_{k}\ \Psi_{K}\right)\\
&\times\exp\left[-\frac{1}{2}\left(\frac{K-K_{\mathrm{o}}}{\Delta K}\right)^{2}-\frac{1}{2}\left(\frac{K'-K_{\mathrm{o}}}{\Delta K}\right)^{2}\right]
\end{align*}
（译注：矩阵元中 ket 的下标原书印作 $K'$；按下文"动量守恒要求 $K'=K+Q$"的叙述，应为 $K$。）

现在，
\begin{equation*}
\int dr\ f(r)e^{iQ\cdot r} = f(Q)
\end{equation*}
$f(Q)$ 有三个重要的性质：(a) $f(Q=0)=(4\pi/3)R_{1}^{3}$，即 $f(r)$ 的"体积"；(b) $f(Q>1/R_{1})\simeq 0$，因而 $f(Q)$ 在 $Q$ 空间中比 $\Delta K$ 陡峭得多地成峰，并且 $f(\Delta K)\simeq 0$；(c) $\sum_{Q}f(Q)=1$。动量守恒要求 $K'=K+Q$，于是性质 (b) 使得指数因子恰好就是
\begin{equation*}
\exp\left[-\left(\frac{K-K_{\mathrm{o}}}{\Delta K}\right)^{2}\right]
\end{equation*}
我们预料，由于 $K$ 与 $K_{\mathrm{o}}$ 非常接近，而矩阵元
\begin{equation*}
M_{K}(Q)
\end{equation*}
没有理由随 $K$ 急速变化，所以我们可以在假定 $M_{K}$ 为常数的情形下完成对 $K$ 的求和：
\begin{equation*}
(6) = e\sum_{Q}\ f(Q)\,M(Q)\ \sum_{k}\ \frac{\text{const.}}{N(\Delta K)^{3}}\ \exp\left[-\left(\frac{K-K_{\mathrm{o}}}{\Delta K}\right)^{2}\right]
\end{equation*}
而这个求和恰恰就是波包的归一化和。于是，电荷表达式 (6) 就变成正好是 $e\sum_{Q}f(Q)M(Q)$。正如我们已经指出的，在 $Q$ 恰好为零时
\begin{equation*}
M(0) = N + 1
\end{equation*}
因此电荷中有一项贡献恰好是
\begin{equation*}
\frac{4\pi R_{1}^{3}}{3}\ e\ (N+1)
\end{equation*}
这一项贡献所代表的，只不过是整个电子气的平均电荷中包含在 $R_{1}$ 球内的那一部分，它当然必须恰好被离子背景电荷所补偿。电荷的这一部分很容易辨认，因为它与波包毫无关系——它会随 $R_{1}$ 无限增大。事实上，可以证明：即使把波包从 $R_{1}$ 球内移走，这一部分也不受影响。因此它纯粹是一个平凡的背景效应。

如果我们令 $M(Q\neq 0)=$ 常数 $M$——这是十分合理的——那么剩下的电荷就是
\begin{equation*}
eM\sum_{Q}\ f(Q) = eM
\end{equation*}
% ---- B.1 收尾（原书 p.112 / PDF p127 顶部，协调者补译；承接上式 eM Σ f(Q) = eM） ----
（由性质 (c)）；于是它与半径 $R_1$ 无关，正是直接与波包自身相联系的那份电荷。正如我们的物理推理早已告诉我们的，它恰好就是 $e/\kappa$。这样，我们就得到了 Kohn 恒等式的一个物理上的——尽管显然不是严格数学意义上的——证明。实际的证明看来只能借助图微扰论（diagrammatic perturbation theory）来进行。应当认识到：整个 $N{+}1$ 体理论的这套机制确实只在微扰论的界限之内被严格证明——尽管恰恰在这个情形下，我们完全有理由相信微扰论实际上收敛，并且准确地代表了物理实在。

Blount 为此补充了一个颇有意思的结果。实质上，我们迄今为止所做的，正相当于 Bloch 电子的零阶单带理论——亦即有效哈密顿量理论
\begin{equation*}
\mathcal{H} = E(\mathbf{k}) + V(\mathbf{R})
\end{equation*}
Blount 指出，只要把 $N{+}1$ 体波函数 $\Psi_{\mathbf{k}}$ 本质上取作
\begin{equation*}
\Psi_{\mathbf{k}}(\mathbf{r}_1,\ \mathbf{r}_2,\ \ldots\,,\mathbf{r}_{N+1})\ =\ \left\{\mathrm{e}^{i\mathbf{k}\cdot\mathbf{r}_1}\ U_{\mathrm{o}}^{\mathbf{k}}\left(\mathbf{r}_1\ldots\mathbf{r}_{N+1}\right)\right\}\ \text{的反对称化}
\end{equation*}
其中 $U_{\mathrm{o}}^{\mathbf{k}}$ 是其 $N{+}1$ 个粒子变量的完全周期函数，就可以得到 $N{+}1$ 体问题的完整理论。于是 $U_{\mathrm{o}}^{\mathbf{k}}$ 的处理方式与通常波函数的 Bloch 部分十分相像，并且由它出发——利用 Kohn 与 Ambegaokar 的恒等式——例如可以导出关系式
\begin{equation*}
V(\mathbf{r}) = \frac{1}{\kappa}\,V(\mathbf{R})\ +\ \text{量级为}\ \frac{\partial V}{\partial \mathbf{R}}\cdot a\ \text{的修正项}
\end{equation*}
然而，这些修正项与真实单电子情形的修正项颇为不同——这是意料之中的，因为现在 $V$ 既能极化原子芯，也能极化所加入电子的波函数。
